\chapter{Thesis Objectives and Deliverables}
\label{chap:objectives}

To fix the problems with older forecasting models and successfully predict when items will run out of stock during the day, our team set the following clear goals for this project:

\begin{enumerate}
    \item \textbf{Set up the Forecasting Problem:} We needed to clearly define the stockout problem mathematically. We wanted our model to look at the last 14 days of history (like sales and discounts) and predict an exact 24-hour breakdown of whether a specific store will have enough stock tomorrow.
    \item \textbf{Test Basic Recurrent Models:} We planned to train basic Vanilla RNN and Vanilla LSTM models first. This helped us see where they struggle and what goes wrong when they try to process messy, real-world retail data.
    \item \textbf{Build Our Own Custom Model:} We aimed to build a specialized neural network. Our final design, the Dual-Stream Inventory Shortcut LSTM, processes the unpredictable demand data separately from the hard inventory facts, and it makes sure the latest inventory updates are not forgotten.
    \item \textbf{Compare Against the Latest Linear Models:} We wanted to test our custom LSTMs against newer, lightweight linear models (like DLinear and Multi-Kernel DLinear) using a very strict testing process to see which approach was actually better.
    \item \textbf{Deliver a Working Solution:} Ultimately, our main goal was to pick one fast, easy-to-understand model (which ended up being the Multi-Kernel DLinear) that gives the most accurate hourly predictions, so it could actually be used in real, fast-paced fulfillment centers.
\end{enumerate}

\section{Problem Formulation}

To set this up, we take historical retail data $X \in \mathbb{R}^{T \times D}$, where $T=14$ days is the time window we look back on, and $D$ is the number of different features we track. Our goal is to train a function $f: X \rightarrow Y$. 

The target output $Y$ is a 24-dimensional vector that tells us exactly if an item is in stock for every hour of the next day ($T+1$):
\begin{equation}
    Y = [y_1, y_2, \dots, y_{24}]^T \quad \text{where} \quad y_i \in \{0, 1\}
\end{equation}
In this setup, if $y_i = 1$, it means the item is out of stock during hour $i$, and if $y_i = 0$, it means there is plenty of inventory.

\section{Dataset Overview}

To hit these goals, our team ran all our experiments using the \textbf{FreshRetailNet-50K} dataset. Our mentor provided this data, and it gives a very realistic picture of how modern retail works.

It includes daily records for specific combinations of cities, stores, and products. Importantly, the dataset mixes direct sales numbers with outside factors that affect how people shop:

\begin{itemize}
    \item \textbf{Sales and Inventory Facts:} Basic numbers like total sales, past stockouts, and daily activity.
    \item \textbf{Promotions:} Information on discounts or store events that make people buy things faster than usual.
    \item \textbf{Outside Events:} National holidays and weather conditions that change overall shopping behavior.
\end{itemize}

By feeding all this information in order, our models learned the difference between normal items selling out over time and items selling out suddenly because of a sale. This dataset paved the way for all the tests we talk about in the next chapters.
